\subsection{LLM-Based Agent Systems}

LLM agents have evolved along a clear trajectory: from prompt-level reasoning loops, to autonomous task decomposition, to domain-specialized production systems. ReAct~\citep{react} and Reflexion~\citep{reflexion} established the basic interaction pattern of alternating reasoning with action and using feedback from prior trajectories to improve future decisions. AutoGPT~\citep{autogpt} then popularized iterative goal decomposition with tool-driven execution, while MetaGPT~\citep{metagpt} made coordination itself a design object by encoding role specialization and SOP-like workflows into multi-agent communication.

Software engineering quickly became the main stress test for these ideas. Devin~\citep{devin} frames coding as long-horizon execution in an integrated environment; SWE-agent~\citep{sweagent} shows that a carefully designed \emph{Agent--Computer Interface} can outperform unrestricted shell access; and OpenHands~\citep{openhands,openhands_sdk} offers an open runtime and SDK for instantiating different agent architectures. Recent product systems push this line further: Claude Code~\citep{claudecode} emphasizes practical coding workflows with explicit context management, OpenAI Codex~\citep{codex} couples coding agents with isolated cloud sandboxes, Manus~\citep{manus} organizes planning, execution, and verification around CodeAct, and OpenClaw~\citep{openclaw_rl} centers interaction around reusable skills.

The common trend is clear: agent systems are becoming more capable and more integrated. Yet most are still optimized either for bounded single-session tasks or for narrow domains such as software engineering. What remains underexplored is how to sustain long-horizon, cross-domain autonomy while keeping context growth under tight control. GenericAgent is motivated by this missing combination of generality and context efficiency.

\subsection{Memory Systems for Agents}

Existing memory work has made long-horizon agents more feasible, but it is much stronger on storage and retrieval than on verification and condensation. MemGPT~\citep{memgpt} treats the context window as ``main memory'' and external storage as ``disk,'' using a paging mechanism to move information across tiers. A-MEM~\citep{amem} instead organizes memory as a dynamically evolving graph of semantically linked atomic notes, enabling richer associative recall.

Production systems adopt more pragmatic variants of the same idea. Claude Code~\citep{claudecode} persists project memory in Markdown files and periodically compacts conversation history, while Manus~\citep{manus} maintains a chronological event trace and external artifacts such as \texttt{todo.md} for plan tracking. These designs extend the effective horizon of the agent, but they remain largely append-only: information is accumulated, summarized, or retrieved, rather than rigorously filtered and distilled into denser validated representations.

This leaves two gaps. First, existing systems provide limited protection against hallucinated or stale information entering persistent memory. Second, they rarely turn accumulated experience into progressively higher-density forms. GenericAgent addresses both through an L0--L3 hierarchical memory and a ``No Execution, No Memory'' rule, so memory updates are grounded in observed outcomes and actively compressed rather than merely accumulated.

\subsection{Self-Evolution and Experience Accumulation}

Research on self-evolving agents asks how experience can be converted into better future behavior, but most methods still stop at textual abstractions. Voyager~\citep{voyager} is the strongest exception: in Minecraft it continually expands a library of verified executable skills that can be recomposed for harder tasks. Beyond such domain-specific lifelong learning, however, the broader literature mainly distills experience into language.

A recent survey on self-evolving agents~\citep{selfevolvesurvey} organizes the space by what evolves, when evolution occurs, and what feedback drives it. Representative methods follow this pattern. EvolveR~\citep{evolver} distills successful trajectories into strategic principles; FLEX~\citep{flex} builds a reflection-based experience library from both successes and failures; AgentEvolver~\citep{agentevolver} stores structured knowledge units in natural language; and the self-evolving lifelong learning framework~\citep{selfevolvinglifelong} stages evolution from exploration to skill learning to knowledge internalization.

The shared limitation is that, outside specialized environments such as Voyager's Minecraft setting, experience is usually retained as notes, principles, or contextual guidance. General-purpose agent systems rarely push distillation all the way into reusable executable procedures. GenericAgent explicitly does so through a three-stage pipeline that turns natural-language experience into structured SOPs and then into code and reusable skills, yielding denser and cheaper representations at inference time.